% Tabla 5: Correspondencia entre la posición nominal de la réplica y la lectura del riel.
% Etiqueta original: tab:posiciones.
% Fuente: capitulos/3-metodologia.tex, lineas 159-178.
% Fragmento extraido automaticamente; requiere el preambulo de main.tex.

\begin{table}[htbp]
\centering\footnotesize
\caption{Correspondencia entre la posición nominal de la réplica y la lectura del riel.}
\label{tab:posiciones}
\begin{tabular}{@{}p{2.8cm}p{3.4cm}p{2.6cm}p{3.4cm}@{}}
\toprule
Posición nominal (cm al vidrio) & Separación de la zona de calibración de B (cm) & Lectura del riel (mm) & Desplazamiento al plano de simetría (mm) \\
\midrule
5  & 0  & 0   & 10, compensado \\
11 & 6  & 60  & 10, compensado \\
17 & 12 & 120 & 10, compensado \\
23 & 18 & 180 & 10, compensado \\
\bottomrule
\end{tabular}

\smallskip
\begin{minipage}{0.9\textwidth}\footnotesize
\emph{Nota.} El plano de simetría de la réplica quedaba 10~mm corrido respecto del apoyo del carrito. Compensamos ese desplazamiento físicamente en el montaje, corriendo el carrito 10~mm para que el plano de simetría quedara en la posición nominal; por eso el procesamiento no aplicó ningún desplazamiento adicional, y el registro de datos consigna un desplazamiento de réplica de 0~mm.
\end{minipage}
\end{table}
